% !TeX root = ../../Book1.tex
% 纲要标题：### 10.2 理想与商环
% 纲要位置：计划纲要.md，第 394 行。

\section{理想与商环}
\label{b1:sec:1002}

在加法群上按一个子群取商，只能保证加法良定义。若还想让乘法通过商映射保留下来，这个子群必须能够吸收环元素的左右乘法，理想由此出现。商环中的零因子与可逆性又将分别引出素理想和极大理想。

% 纲要内容（仅作写作计划，尚非教材正文）：
%
% * 左理想、右理想与双边理想
% * 商环及同构定理
% * 素理想与极大理想的首次出现
%

\subsection{左理想、右理想与双边理想}
\label{b1:subsec:1002:01}
设 \(R\) 为环，先取加法子群 \(I\leq(R,+)\)，希望陪集上的乘法由代表元的乘法给出。若把第一项的代表 \(a\) 换成 \(a+\ell\)，其中 \(\ell\in I\)，所得乘积与原乘积的差是
\[
(a+\ell)b-ab=\ell b.
\]
要使输出陪集不变，必须有 \(\ell b\in I\)。若换的是第二项的代表，同样由 \(a(b+\ell)-ab=a\ell\) 得到 \(a\ell\in I\)。所以 \(I\) 既要吸收右乘，也要吸收左乘。与\cref{b1:thm:quotient-group}中由换代表元引出正规性一样，这里的条件也来自商运算的良定义要求。

\begin{definition}\label{b1:def:ring-ideal}
设 \(R\) 为环，\(I\subseteq R\) 是包含零元且对减法封闭的子集。若对每个 \(r\in R,a\in I\) 都有 \(ra\in I\)，称 \(I\) 为\term[zuo-lixiang]{左理想}[left ideal]；若总有 \(ar\in I\)，称为\term[you-lixiang]{右理想}[right ideal]；同时满足两者时称为\term[shuangbian-lixiang]{双边理想}[two-sided ideal]。未加限定的“理想”均指双边理想。若 \(I\ne R\)，称为真理想。
\end{definition}
理想不要求含环的单位；事实上理想一旦含有单位元或任一单位 \(u\)，就包含 \(u^{-1}u=1\)，继而包含每个 \(r=r1\)，所以等于 \(R\)。因此，一个非零环的真理想不是本书约定下的含同一单位的子环。这个限定不排除它自身的运算：理想仍对加法、减法和乘法封闭，可以作为不要求单位元的环来研究，有时也具有自己的单位元；这里排除的是它含有母环单位 \(1_R\)。

\ProseParagraphBreak
在交换环中左右条件等价。\(n\mathbb Z\) 是 \(\mathbb Z\) 的理想；反过来，整数环的每个理想首先是加法子群，由\cref{b1:lem:integer-subgroups}都为这种形式。

\ProseParagraphBreak
非交换情形确须区分左右。设 \(k\) 为域，在 \(M_2(k)\) 中令
\[
I_\ell=\left\{\,\begin{pmatrix}a&0\\b&0\end{pmatrix}
\;\middle|\;a,b\in k\,\right\},\qquad
I_r=\left\{\,\begin{pmatrix}a&b\\0&0\end{pmatrix}
\;\middle|\;a,b\in k\,\right\}.
\]
它们分别由第二列为零、第二行为零的矩阵组成，都对减法封闭。对任意 \(a,b,r,s,t,u\in k\)，
\[
\begin{aligned}
\begin{pmatrix}r&s\\t&u\end{pmatrix}
\begin{pmatrix}a&0\\b&0\end{pmatrix}
&=\begin{pmatrix}ra+sb&0\\ta+ub&0\end{pmatrix},\\
\begin{pmatrix}a&b\\0&0\end{pmatrix}
\begin{pmatrix}r&s\\t&u\end{pmatrix}
&=\begin{pmatrix}ar+bt&as+bu\\0&0\end{pmatrix}.
\end{aligned}
\]
第一式说明 \(I_\ell\) 吸收左乘，第二式说明 \(I_r\) 吸收右乘。但 \(E_{11}\in I_\ell\) 而 \(E_{11}E_{12}=E_{12}\notin I_\ell\)，所以 \(I_\ell\) 不是右理想；\(E_{11}\in I_r\) 而 \(E_{21}E_{11}=E_{21}\notin I_r\)，所以 \(I_r\) 不是左理想。

\ProseParagraphBreak
若 \(f:R\rightarrow S\) 为环同态，其加法核 \(\ker f=\{a\in R:f(a)=0\}\) 为双边理想：对 \(a\in\ker f\)，有 \(f(ra)=f(r)f(a)=0=f(ar)\)。这里保持零元的条件由加法保证，吸收乘法的条件来自乘法保持性。

\subsection{商环及同构定理}
\label{b1:subsec:1002:02}
\begin{proposition}\label{b1:prop:quotient-ring}
设 \(R\) 为环，\(L\leq(R,+)\) 为加法子群。在加法陪集集合 \(R/L\) 上按
\[
(a+L)+(b+L)=(a+b)+L,\qquad
(a+L)(b+L)=ab+L
\]
规定运算，则乘法与代表元无关当且仅当 \(L\) 是 \(R\) 的双边理想。在此情形下，\(R/L\) 构成\term[shanghuan]{商环}[quotient ring]，单位为 \(1+L\)；自然映射 \(q:R\rightarrow R/L\) 是满环同态，核为 \(L\)。
\end{proposition}
\begin{proof}
加法良定义来自加法群的商构造。若 \(L\) 是双边理想，而 \(a'=a+\ell,b'=b+m\)，其中 \(\ell,m\in L\)，则
\[
a'b'-ab=am+\ell b+\ell m\in L.
\]
其中 \(am\in L\) 用左吸收，因为 \(m\in L\)、\(a\in R\)；\(\ell b\in L\) 用右吸收，因为 \(\ell\in L\)、\(b\in R\)；\(\ell m\) 则用任一侧吸收都可得到。故乘法与代表元无关。结合律和分配律由代表元上的对应等式下降得到，\(1+L\) 在两侧乘任一陪集都保持它。自然映射按定义保持两种运算与单位，陪集都有代表元，且 \(a+L=L\) 恰好等价于 \(a\in L\)。

反过来，设所给陪集乘法良定义。对任意 \(r\in R\) 与 \(\ell\in L\)，由 \(\ell+L=0+L\) 得
\[
\begin{aligned}
r\ell+L&=(r+L)(\ell+L)=(r+L)(0+L)=L,\\
\ell r+L&=(\ell+L)(r+L)=(0+L)(r+L)=L.
\end{aligned}
\]
于是 \(r\ell,\ell r\in L\)。\(L\) 已是加法子群，故为双边理想。
\end{proof}
\symidx{\(R/I\)：环关于双边理想的商环}

把环同态 \(f:R\to S\) 先看作加法群同态，\cref{b1:thm:first-isomorphism}已经给出加法群同构 \((R,+)/\ker f\cong\operatorname{im}f\)。环论需要补的检查是：核吸收左右乘法，使商乘法有意义；诱导映射除保持加法外，还保持乘法和单位。前面已经证明核是双边理想，像则对两种运算封闭并含 \(1_S=f(1_R)\)，是 \(S\) 的子环。因而可以沿用同一个陪集映射得到下面的环同构。

\begin{theorem}[环的第一同构定理]\label{b1:thm:ring-first-iso}
设 \(R,S\) 为环。若 \(f:R\rightarrow S\) 是环同态，则
\[
\overline f:R/\ker f\longrightarrow\operatorname{im}f,
\qquad a+\ker f\longmapsto f(a)
\]
为环同构。更一般地，若 \(I\) 为 \(R\) 的双边理想且 \(I\subseteq\ker f\)，并以 \(q:R\rightarrow R/I\) 表示自然映射，则存在唯一环同态 \(g:R/I\rightarrow S\) 满足 \(g\circ q=f\)。
\end{theorem}
\begin{proof}
若 \(a-a'\in I\subseteq\ker f\)，则 \(f(a)=f(a')\)，所以 \(g(a+I)=f(a)\) 良定义。加法相容性由加法群的商构造给出，乘法与单位则分别满足
\[
\begin{aligned}
g\bigl((a+I)(b+I)\bigr)&=f(ab)=f(a)f(b),\\
g(1_R+I)&=f(1_R)=1_S.
\end{aligned}
\]
所以 \(g\) 是保单位的环同态。由定义有 \(g\circ q=f\)；由于 \(q\) 满射，该公式也是满足此等式的同态的唯一可能取值。取 \(I=\ker f\)，并将 \(g\) 的目标限制为其实际像 \(\operatorname{im}f\)，便得到定理中的 \(\overline f\)。它按像的定义满射；若 \(\overline f(a+\ker f)=0\)，则 \(a\in\ker f\)，所以其核只有零陪集，因而 \(\overline f\) 也是单射。
\end{proof}
条件 \(I\subseteq\ker f\) 使下面的三角形交换，虚线表示唯一满足 \(g\circ q=f\) 的环同态。
\[
\begin{tikzcd}[column sep=3em,row sep=2.5em]
R\arrow[r,"q"]\arrow[dr,"f"']&R/I\arrow[d,dashed,"\exists!\,g"]\\
&S
\end{tikzcd}
\]
这里的泛性质告诉我们：要从商环定义映射，先在原环定义，并检查需要抹去的理想确实映成零。商环不只是一个同构类型，还连同指定的商映射一起发挥作用。

\begin{proposition}[理想对应与商的同构]\label{b1:prop:ring-ideal-correspondence}
设 \(R\) 为环，\(I\) 为 \(R\) 的双边理想，\(q:R\rightarrow R/I\) 为自然映射。则：
\begin{enumerate}
\item 理想对应：\(q\) 给出包含 \(I\) 的 \(R\) 的双边理想与 \(R/I\) 的双边理想之间的保序双射。令 \(J\) 遍历前一类理想，\(L\) 遍历后一类理想，双向公式为
\begin{equation}\label{b1:eq:ring-ideal-correspondence}
J\longmapsto J/I,\qquad L\longmapsto q^{-1}(L).
\end{equation}
\item 环的第三同构定理：若 \(J\) 为 \(R\) 的双边理想且 \(I\subseteq J\)，则
\begin{equation}\label{b1:eq:ring-third-iso}
(R/I)/(J/I)\cong R/J.
\end{equation}
\item 环的第二同构定理：若 \(A\subseteq R\) 为含同一单位的子环，并记 \(A+I=\{a+i:a\in A,\ i\in I\}\)，则 \(A+I\) 是 \(R\) 的子环，\(I\) 是 \(A+I\) 的双边理想，且
\begin{equation}\label{b1:eq:ring-second-iso}
A/(A\cap I)\cong(A+I)/I.
\end{equation}
\end{enumerate}
\end{proposition}
\begin{proof}
\begin{enumerate}
\item 设 \(L\) 是 \(R/I\) 的双边理想。其原像 \(q^{-1}(L)\) 对减法及左右乘法封闭，且包含 \(I\)。含 \(I\) 的理想 \(J\) 的像 \(q(J)=J/I\) 对 \(R/I\) 中元素的左右乘法封闭，因为这些元素都可提升到 \(R\)。有 \(q^{-1}(q(J))=J\)：若 \(q(a)=q(j)\)，则 \(a-j\in I\subseteq J\)。又 \(q(q^{-1}(L))=L\)，因为 \(q\) 满射。两种操作均保持包含关系，得到\cref{b1:eq:ring-ideal-correspondence}。
\item 因 \(I\subseteq J\)，映射 \(R/I\rightarrow R/J\)、\(a+I\mapsto a+J\) 良定义且满射，核为 \(J/I\)。对该映射应用\cref{b1:thm:ring-first-iso}，即得\cref{b1:eq:ring-third-iso}。
\item 对 \(a,b\in A\) 及 \(i,j\in I\)，两元素 \(a+i,b+j\) 的差属于 \(A+I\)，乘积
\[
(a+i)(b+j)=ab+aj+ib+ij\in A+I.
\]
假设 \(A\) 含母环的同一个单位正用于这里：\(1_R\in A\subseteq A+I\)，配合上述封闭性，使 \(A+I\) 成为本书约定下的子环；只有加法与乘法封闭还不足以保证这一点。\(I\subseteq A+I\) 且吸收其左右乘法，因此是它的双边理想。\(A\cap I\) 同样是 \(A\) 的双边理想。映射 \(A\rightarrow(A+I)/I\)、\(a\mapsto a+I\) 保持单位且满射，因为 \(1_R\mapsto1_R+I\)、\(a+i+I=a+I\)，其核为 \(A\cap I\)。应用\cref{b1:thm:ring-first-iso}，即得\cref{b1:eq:ring-second-iso}。
\end{enumerate}
\end{proof}
上面的理想对应和两个商同构分别来自不同的映射。表中 \(I\subseteq J\) 是 \(R\) 的双边理想，\(A\) 是含 \(1_R\) 的子环；把出发映射及其核写出，就能看清每条公式抹去了什么。

\begin{center}
\begin{tabularx}{\linewidth}{@{}l>{\raggedright\arraybackslash}Xl@{}}
\toprule
结论 & 出发映射 & 核 \\
\midrule
理想对应 & \(q:R\to R/I\)，\(a\mapsto a+I\) & \(I\) \\
第三同构 & \(R/I\to R/J\)，\(a+I\mapsto a+J\) & \(J/I\) \\
第二同构 & \(A\to(A+I)/I\)，\(a\mapsto a+I\) & \(A\cap I\) \\
\bottomrule
\end{tabularx}
\end{center}

环同态的像未必是目标环的理想，例如 \(\mathbb Z\hookrightarrow\mathbb Q\) 的像不是 \(\mathbb Q\) 的理想。上述理想对应使用商映射的满射性，不能省去这个条件。

\subsection{素理想与极大理想}
\label{b1:subsec:1002:03}
本小节只讨论交换环。非交换环的素性与极大理想需要区分不同版本，不把下面的乘积判别直接当作其一般定义。

\ProseParagraphBreak
取真理想 \(I\) 后，商环 \(R/I\) 非零。接下来有两个不同的问题：什么时候它仍像整数环一样，两个非零元素的乘积不会为零？什么时候它像域一样，每个非零元素都能取逆？下面的素理想与极大理想将分别刻画这两种商环性质。

\begin{definition}\label{b1:def:prime-maximal-ideal}
设 \(R\) 为交换环，\(\mathfrak p\subsetneq R\) 为真理想。若对任意 \(a,b\in R\)，\(ab\in\mathfrak p\) 都蕴含 \(a\in\mathfrak p\) 或 \(b\in\mathfrak p\)，则称 \(\mathfrak p\) 为\term[su-lixiang]{素理想}[prime ideal]。再设 \(\mathfrak m\subsetneq R\) 为真理想。若满足 \(\mathfrak m\subseteq J\subseteq R\) 的理想只有 \(J=\mathfrak m\) 与 \(J=R\)，则称 \(\mathfrak m\) 为\term[jida-lixiang]{极大理想}[maximal ideal]。
\end{definition}
“真”不能漏掉；否则整个环会满足乘积条件，却不能对应非零整环。\((0)\) 是整环中的素理想，但整数环中它不极大，因为 \((0)\subsetneq2\mathbb Z\subsetneq\mathbb Z\)。

\begin{proposition}\label{b1:prop:prime-max-quotient}
交换环 \(R\) 的理想 \(I\) 为素理想，当且仅当 \(R/I\) 是整环；为极大理想，当且仅当 \(R/I\) 是域。因此每个极大理想都是素理想。
\end{proposition}
\begin{proof}
\(I\ne R\) 恰好保证商环非零；在商环中，\((a+I)(b+I)=0\) 恰好等价于 \(ab\in I\)，故第一项直接由定义得到。

若 \(I\) 极大且 \(a\notin I\)，要使非零类 \(a+I\) 有逆，就要找到 \(r\in R\) 使 \(ra-1\in I\)。这等价于 \(1\in I+Ra\)，所以考虑包含 \(I\) 并加入 \(a\) 的集合
\[
I+Ra=\{i+ra:i\in I,\ r\in R\}.
\]
此集合对减法封闭；由于 \(R\) 交换，对 \(s\in R\) 还有 \(s(i+ra)=si+(sr)a\in I+Ra\)，故它是理想。它含 \(I\) 及 \(a\)，而 \(a\notin I\)，所以严格包含 \(I\)，由极大性知 \(I+Ra=R\)。写 \(1=i+ra\)，便有 \((r+I)(a+I)=1+I\)，所以商中每个非零元素可逆。反过来，若 \(R/I\) 为域，而理想 \(J\) 严格包含 \(I\)，取 \(a\in J\setminus I\)；商中 \(a+I\) 的逆有代表 \(r\)，于是 \(1-ra\in I\subseteq J\)，给出 \(1\in J\)，故 \(J=R\)。最后，域无非零零因子，所以极大理想为素理想。
\end{proof}
素理想不一定极大，正如整环不一定是域：\(\mathbb Z/(0)\cong\mathbb Z\) 是整环，却不是域，所以 \((0)\) 是 \(\mathbb Z\) 的素理想而不是极大理想。这与前面的严格包含链 \((0)\subsetneq2\mathbb Z\subsetneq\mathbb Z\) 给出同一个反例。

\ProseParagraphBreak
整数环的理想为 \(n\mathbb Z\)。当 \(p\) 是素数，任一 \(a\notin p\mathbb Z\) 满足 \(\gcd(a,p)=1\)，整数 Bézout 等式给出 \(ua+vp=1\)，故模 \(p\) 的非零类均有逆。所得域记为 \(\mathbb F_p=\mathbb Z/p\mathbb Z\)。正合数 \(n=ab\)、\(1<a,b<n\) 则产生两个非零类之积为零。因此 \(\mathbb Z\) 的素理想恰为 \((0)\) 及 \(p\mathbb Z\)，极大理想恰为后者。
\symidx{\(\mathbb F_p\)：含 \(p\) 个元素的素数域}

\ProseParagraphBreak
极大理想的存在性寻找的是包含偏序中的极大元，即不能再严格扩大的真理想，并非包含所有真理想的最大元。例如 \(2\mathbb Z\) 与 \(3\mathbb Z\) 都是极大理想，却互不包含；没有哪一个同时包含这两个真理想并仍为真理想。下面用 Zorn 引理证明所需极大元存在。

\begin{proposition}\label{b1:prop:maximal-ideal-exists}
非零交换环中的每个真理想都包含在某个极大理想中。
\end{proposition}
\begin{proof}
固定真理想 \(I\)，考虑包含 \(I\) 的真理想按包含关系构成的偏序集。它非空。非空链的并对减法封闭，因为任意两个元素都落在链上某个共同的理想中；它也吸收乘法。这个并不含 \(1\)，否则链中的某个真理想已含 \(1\)，矛盾。因而链的并仍是此偏序集的上界；空链可用 \(I\) 作上界。应用 Zorn 引理（\cref{b1:thm:A-zorn}），即“每条链有上界的非空偏序集有极大元”，得到 \(\mathfrak m\)。任何严格包含它的真理想仍包含 \(I\)，违背极大性，故 \(\mathfrak m\) 是极大理想。
\end{proof}
这里明确使用选择公理的 Zorn 形式，不给出寻找极大理想的算法。在具体环中仍应尽可能写出商映射，例如对交换环 \(R\)，常数项映射 \(R[x]\rightarrow R\) 满射，核为 \(xR[x]\)。\cref{b1:thm:ring-first-iso}给出 \(R[x]/xR[x]\cong R\)，所以当 \(R\) 为整环或域时，这个理想分别为素理想或极大理想。
